% =====================================================
% Modern CV — styling commands
% =====================================================

% --- 1. Accent colour ---
\definecolor{cvaccent}{HTML}{006795}      % MidnightBlue
\definecolor{cvmuted}{HTML}{6B7280}       % grey for dates / secondary text

% --- 2. Inline highlighted phrase (bold, MidnightBlue, TitleCase) ---
\newcommand{\txthl}[1]{\textcolor{cvaccent}{\textbf{\titlecap{#1}}}}

% --- 3. Top-level CV section heading ----------------------------------
% NB: \noindent before the \rule line is essential — without it, the
% paragraph starts with \parindent (15pt) and the rule overshoots the
% margin by exactly that amount.
% The heading block itself consumes ~2.5\baselineskip (vspace + \Large
% title line + rule + vspace). What follows immediately (e.g.
% \cvsmallheading) does its OWN separate \needspace check (5\baselineskip)
% that is only evaluated AFTER the heading has already been placed --
% so the two checks don't automatically compose. To guarantee the
% heading is never left alone at the bottom of a page, this value must
% cover the heading's own height PLUS whatever the very next \needspace
% check downstream demands (2.5 + 5 = 7.5, so 10 gives comfortable
% margin). 4, then 7\baselineskip were tried first and both still left
% "Clients" (CV_Consulting.tex) orphaned -- confirmed by rendering the
% compiled page.
\newcommand{\cvsection}[1]{%
  \needspace{10\baselineskip}%
  \par
  \vspace{0.9\baselineskip}%
  \noindent
  {\color{cvaccent}{\Large\bfseries \titlecap{#1}}\par}%
  \vspace{-0.05\baselineskip}%
  \noindent{\color{cvaccent}\rule{\linewidth}{0.4pt}\par}%
  \vspace{0.4\baselineskip}%
}

% --- 3b. Same as cvsection but without titlecap (safe for headings that
%        contain commands like \vspace, \\ or punctuation like '–' that
%        would confuse \titlecap). Use for legacy already-cased titles. ---
\newcommand{\cvsectionraw}[1]{%
  \needspace{10\baselineskip}%
  \par
  \vspace{0.9\baselineskip}%
  \noindent
  {\color{cvaccent}{\Large\bfseries #1}\par}%
  \vspace{-0.05\baselineskip}%
  \noindent{\color{cvaccent}\rule{\linewidth}{0.4pt}\par}%
  \vspace{0.4\baselineskip}%
}

% --- 4. Block-style sub-section (smaller than cvsection) --------------
\newcommand{\cvsubsection}[1]{%
  \needspace{5\baselineskip}%
  \par
  \vspace{0.5\baselineskip}%
  \noindent{\color{cvaccent}\bfseries \titlecap{#1}}\par
  \vspace{0.15\baselineskip}%
}

% --- 5. Inline lead-in heading (runs into following text) -------------
\newcommand{\somespace}{\vspace{0.55em}}
\newcommand{\cvsmallheading}[1]{\needspace{5\baselineskip}\par\somespace\noindent\txthl{#1}\hspace{0.3em}\ignorespaces}

% --- 5b. Same as cvsmallheading but without titlecap (safe for headings
%        that contain a literal '–' or other punctuation that breaks
%        \titlecap's word-scanning under pdfLaTeX -- same rationale as
%        \cvsectionraw above). Use for already-cased text. -------------
\newcommand{\txthlraw}[1]{\textcolor{cvaccent}{\textbf{#1}}}
\newcommand{\cvsmallheadingraw}[1]{\needspace{5\baselineskip}\par\somespace\noindent\txthlraw{#1}\hspace{0.3em}\ignorespaces}

% --- 6. CV entry: date | role | details -------------------------------
\newlength{\datebox}
\setlength{\datebox}{1.7cm}
\newcommand{\cventry}[3]{%
  \noindent
  \parbox[t]{\datebox}{\raggedright\small\color{cvmuted} #1}%
  \hspace{0.5em}%
  \parbox[t]{\dimexpr\linewidth-\datebox-0.6em\relax}{%
    \textbf{\leavevmode #2}\\[0.15\baselineskip]%
    \small #3%
  }%
  \par\vspace{0.35\baselineskip}%
}

% --- 6b. Compact dated entry: date column | description text.
%        Uses hanging-indent layout so (a) the year sits on exactly
%        the same baseline as the first line of text, and (b) all
%        continuation lines start at the same horizontal position
%        regardless of whether the date is a single year or a range.
%        \cvdatesep is the gap between the date box and the text. ------
\newlength{\cvdatesep}
\setlength{\cvdatesep}{0.45em}
\newcommand{\cvitem}[2]{%
  \par\noindent
  \hangindent=\dimexpr\datebox+\cvdatesep\relax
  \hangafter=1
  \makebox[\datebox][l]{\small\color{cvmuted} #1}%
  \hspace{\cvdatesep}%
  {\small #2}%
  \par\vspace{0.15\baselineskip}%
}

% --- 6c. \cvitemwide is now identical to \cvitem (kept for compatibility
%        with existing module files — the old wider-box variant is no
%        longer needed since \datebox is set consistently). -------------
\newlength{\datebigbox}
\setlength{\datebigbox}{2.2cm}
\newcommand{\cvitemwide}[2]{\cvitem{#1}{#2}}

% --- 7. Section band (used as Part divider) ---------------------------
\NewDocumentCommand{\sectionband}{ O{cvaccent} O{white} m }{%
  \par\medskip
  \noindent
  \begingroup
    \setlength{\fboxsep}{0.7ex}%
    \colorbox{#1}{%
      \parbox{\dimexpr\linewidth-2\fboxsep\relax}{%
        \raggedright\bfseries\Large \textcolor{#2}{#3}%
      }%
    }%
  \endgroup
  \par\smallskip
}

% --- 8. Author name + journal helpers (used in Publications module) ---
\providecommand{\rt}{{\bfseries R.~Trotta}\xspace}
\providecommand{\etal}{{\itshape et al.}\xspace}

% =====================================================
% Compatibility layer for existing module files
% (lets the original CV_Academia.tex, CV_Publications.tex, etc.
% compile under the new style without rewriting them)
% =====================================================

% Map legacy headings onto new style. Legacy headings are already in title
% case and sometimes contain inline commands (\vspace, \\) so we route them
% through \cvsectionraw, which does NOT apply \titlecap.
\providecommand{\maintitle}[1]{\needspace{6\baselineskip}\sectionband{#1}}
% (\cvsectionraw already does its own \needspace, so no need to repeat it here.)
\providecommand{\tlnew}[1]{\cvsectionraw{#1}}
\providecommand{\tlnewnospace}[1]{\cvsectionraw{#1}}
\providecommand{\subt}[1]{\needspace{3\baselineskip}\par\medskip\noindent{\bfseries\itshape\color{cvaccent}#1}\par\smallskip}
\providecommand{\sectitle}[1]{\cvsectionraw{#1}}

% Column widths used by original tabular layouts.
% colthreewidth was set to 14cm in the legacy modules, which sits ~15pt
% past the available text width with 2cm side margins under \small in
% Part II — the source of all the "Overfull \hbox (15.0pt too wide)"
% warnings. 13.4cm fits and leaves a small breathing margin.
\providecommand{\colonewidth}{1.4cm}
\providecommand{\colthreewidth}{13.4cm}
\providecommand{\colthreewidthshort}{13.4cm}
\providecommand{\colonewidthlong}{1.9cm}

% Belt and braces: allow LaTeX to stretch interword space a little to
% absorb stubborn long URLs / titles before flagging an overfull box.
\setlength{\emergencystretch}{3em}

% Misc spacing & visual macros
\providecommand{\spa}{\rule{0pt}{11pt}}
\providecommand{\extraspace}{\rule{0pt}{13pt}}
\providecommand{\linea}{$\blacktriangleright$~}
\providecommand{\so}[1]{#1}                     % soul letter-spacing → no-op
\providecommand{\ul}[1]{\underline{#1}}         % soul underline      → \underline
\providecommand{\rtcomment}[1]{}                % silence draft comments
\providecommand{\testimonial}[2]{{\itshape #1}\par\hfill #2\par}
\providecommand{\cites}[1]{{\bfseries Citations: #1}}
\providecommand{\mycite}[1]{{\small [#1]}}

% =====================================================
% Entry-counting infrastructure
% =====================================================
\makeatletter

% Active counter name (empty = off)
\newcommand{\cv@currentcounter}{}
\newcommand{\setcvcounter}[1]{\renewcommand{\cv@currentcounter}{#1}}
\newcommand{\clearcvcounter}{\renewcommand{\cv@currentcounter}{}}
\newcommand{\stepcvcurrent}{%
  \ifx\cv@currentcounter\@empty\else\stepcounter{\cv@currentcounter}\fi}

% \cvcount{year}{text}: like \cvitem but increments the active counter
\newcommand{\cvcount}[2]{\stepcvcurrent\cvitem{#1}{#2}}

% \cvsep: elegant separator between merged entries (also increments counter).
% Uses a forced line-break with small extra gap, then indents to the text
% column and draws a small accent-coloured marker.
\newlength{\cvsepvskip}
\setlength{\cvsepvskip}{0.22\baselineskip}
\newcommand{\cvsep}{%
  \stepcvcurrent
  \unskip\\[\cvsepvskip]%
  \hspace*{0.5cm}%
  {\scriptsize\color{cvaccent}$\triangleright$}\enspace
  \ignorespaces}

% Write a counter's current value to the .aux file so Part I can \ref it.
\newcommand{\labelcvcount}[1]{%
  \protected@write\@auxout{}{%
    \string\newlabel{cvcount@#1}{{\arabic{#1}}{}{}{}{}}}%
}

\makeatother

% Reference a saved counter value (resolves after second LaTeX pass).
\newcommand{\refcvcount}[1]{\ref{cvcount@#1}}

% ── counters ──────────────────────────────────────────────────────────────
% Scientific publications
\newcounter{nRefPhysicsPapers}
\newcounter{nRefMLPapers}
\newcounter{nRefSCPapers}
\newcounter{nAllRefPapers}
\newcounter{nBookChapters}
\newcounter{nWhitePapers}
\newcounter{nProceedings}
% PE publications
\newcounter{nPEBooks}
\newcounter{nPEBookChapters}
\newcounter{nPEOtherPubs}
\newcounter{nAllPEPubs}
% PE activities
\newcounter{nPublicLectures}
\newcounter{nSCTeaching}
\newcounter{nCoursesPublic}
\newcounter{nArtProjects}
\newcounter{nBroadcastMedia}
\newcounter{nPrintCoverage}
% Academia
\newcounter{nInvitedTalks}
\newcounter{nContribTalks}
\newcounter{nOrgMeetings}

% =====================================================
% Grant value totals (used by \grant / \grantE in CV_Grants.tex)
% =====================================================
% Running accumulators (plain decimal strings, in currency-k units),
% reset by \grantsubtotal after each subsection. Exact decimal
% arithmetic via xfp (\fpeval) -- no floating point/rounding artefacts.
% "grand" accumulators are never reset -- they run over the whole
% grant list and feed the overall total at the top (via the .aux file,
% same forward-reference trick as \pageref{LastPage} elsewhere in the CV).
\newcommand\gGBPtotal{0}
\newcommand\gGBPpi{0}
\newcommand\gEURtotal{0}
\newcommand\gEURpi{0}
\newcommand\gGBPgrand{0}
\newcommand\gGBPgrandpi{0}
\newcommand\gEURgrand{0}
\newcommand\gEURgrandpi{0}

\newcommand\gAddGBP[1]{\xdef\gGBPtotal{\fpeval{\gGBPtotal+(#1)}}\xdef\gGBPgrand{\fpeval{\gGBPgrand+(#1)}}}
\newcommand\gAddGBPpi[1]{\xdef\gGBPpi{\fpeval{\gGBPpi+(#1)}}\xdef\gGBPgrandpi{\fpeval{\gGBPgrandpi+(#1)}}}
\newcommand\gAddEUR[1]{\xdef\gEURtotal{\fpeval{\gEURtotal+(#1)}}\xdef\gEURgrand{\fpeval{\gEURgrand+(#1)}}}
\newcommand\gAddEURpi[1]{\xdef\gEURpi{\fpeval{\gEURpi+(#1)}}\xdef\gEURgrandpi{\fpeval{\gEURgrandpi+(#1)}}}

% Format a plain number (currency-k) as "Xk" or, above 1000, "XM",
% rounded to 2 decimal places after the point in the displayed unit
% (trailing zeros are dropped by xfp, e.g. 2.00 -> "2").
\newcommand\gFmtMoney[1]{%
  \ifnum\fpeval{(#1)>=1000}=1
    \fpeval{round((#1)/1000,2)}M%
  \else
    \fpeval{round(#1,2)}k%
  \fi
}

% Bold a display string #2 if its numeric value #1 (currency-k) is >=100k
% -- used to auto-highlight grants above 100k in any currency, without
% touching the surrounding item text.
\newcommand\gBoldIfBig[2]{\ifnum\fpeval{(#1)>=100}=1\textbf{#2}\else#2\fi}

% Print a GBP/EUR total + PI-only sub-total pair (shared by \grantsubtotal
% and \grantgrandtotal). #1=GBP total #2=GBP PI #3=EUR total #4=EUR PI.
\newcommand\gPrintTotals[4]{%
  \ifnum\fpeval{#1>0}=1
    \pounds\gFmtMoney{#1}%
    \ifnum\fpeval{#2>0}=1
      \space(of which \pounds\gFmtMoney{#2} as PI)%
    \fi
  \fi
  \ifnum\fpeval{#1>0}=1\ifnum\fpeval{#3>0}=1\relax\space\textbullet\space\fi\fi
  \ifnum\fpeval{#3>0}=1
    \euro\gFmtMoney{#3}%
    \ifnum\fpeval{#4>0}=1
      \space(of which \euro\gFmtMoney{#4} as PI)%
    \fi
  \fi
}

% \grantsubtotal: print the GBP/EUR totals (and PI-only sub-totals)
% accumulated since the previous \grantsubtotal (or the start of the
% grant list), then reset the accumulators for the next subsection.
\newcommand\grantsubtotal{%
  \par\vspace{0.15\baselineskip}%
  \noindent{\small\itshape
    \textbf{Subtotal:}\space
    \gPrintTotals{\gGBPtotal}{\gGBPpi}{\gEURtotal}{\gEURpi}%
  }\par
  \xdef\gGBPtotal{0}\xdef\gGBPpi{0}%
  \xdef\gEURtotal{0}\xdef\gEURpi{0}%
}

% \grantgrandtotal: print the overall GBP/EUR totals (+ PI sub-totals)
% across the whole grant list. Relies on \gWriteGrantGrandTotal (called
% at the end of the list) having written the values on a previous run.
% Read back via \getrefnumber (from refcount, loaded by hyperref) rather
% than \ref: plain \ref would come back wrapped in hyperref's hyperlink
% markup, which \fpeval cannot parse as a number; \getrefnumber returns
% the bare value (and "0" if it isn't known yet, on the very first run).
\newcommand\grantgrandtotal{%
  \par\noindent{\small\itshape
    \textbf{Overall total:}\space
    \gPrintTotals{\getrefnumber{grantgrand@GBP}}{\getrefnumber{grantgrand@GBPpi}}{\getrefnumber{grantgrand@EUR}}{\getrefnumber{grantgrand@EURpi}}%
  }\par
}

\makeatletter
\newcommand\gWriteGrantGrandTotal{%
  \protected@write\@auxout{}{\string\newlabel{grantgrand@GBP}{{\gGBPgrand}{}{}{}{}}}%
  \protected@write\@auxout{}{\string\newlabel{grantgrand@GBPpi}{{\gGBPgrandpi}{}{}{}{}}}%
  \protected@write\@auxout{}{\string\newlabel{grantgrand@EUR}{{\gEURgrand}{}{}{}{}}}%
  \protected@write\@auxout{}{\string\newlabel{grantgrand@EURpi}{{\gEURgrandpi}{}{}{}{}}}%
}
\makeatother

%%%%%%%%%%%%%%%%%%%%%%%%
% Journals' names 
%%%%%%%%%%%%%%%%%%%%%%%%
\newcommand{\mnras}{{\em Mon.\ Not.\ R.~Astron.~Soc.~}}
\newcommand{\mnraslett}{{\em Mon.~Not.~R.~Astron.~Soc.~Lett.}}
\newcommand{\PRL}{{\em Phys.~Rev.~Lett.~}}
\newcommand{\PRD}{{\em Phys.~Rev.~D~}}
\newcommand{\PRB}{{\em Phys.~Rev.~B~}}
\newcommand{\JHEP}{{\em JHEP~}}
\newcommand{\JCAP}{{\em JCAP~}}
\newcommand{\ApJ}{{\em ApJ~}}
\newcommand{\ApJL}{{\em ApJL~}}
\newcommand{\PLB}{{\em Phys. Lett. B~}}
\newcommand{\NAR}{{\em New~Astron.~Rev.}}
